% Part: normal-modal-logic
% Chapter: syntax-and-semantics
% Section: relational-models

\documentclass[../../../include/open-logic-section]{subfiles}

\begin{document}

\olfileid{nml}{syn}{rel}

\olsection{Model Relasional}

Konsep dhasar semantik kanggo logika modal normal yaiku
\emph{model relasional}. Model iki dumadi saka himpunan donya kang
disambungake dening ``relasi aksesibilitas'' biner, sarta penetapan
kang nemtokake !!{propositional variable}s endi kang dianggep
``bener'' ing donya endi.

\begin{defn}
  Sawijining \emph{model} kanggo basa modal dhasar iku tripel $\mModel{M}
  = \tuple{W, R, V}$, kanthi
  \begin{enumerate}
  \item $W$ iku himpunan ``donya'' kang ora kosong,
  \item $R$ iku relasi aksesibilitas biner ing~$W$, lan
  \item $V$ iku fungsi kang nemtokake kanggo saben !!{propositional
    variable}~$p$ sawijining himpunan $V(p)$ saka donya kang mungkin.
  \end{enumerate}
  Yen $Rww'$ lumaku, kita kandha yen $w'$ \emph{bisa diakses
    saka}~$w$. Yen $w \in V(p)$, kita kandha yen $p$ \emph{bener ing}~$w$.
\end{defn}

Kauntungan utama semantik relasional yaiku model bisa
digambarake nganggo diagram prasaja, kaya ing \olref{fig:simple}.
Donya digambarake minangka simpul, lan donya~$w'$ bisa diakses
saka~$w$ persis yen ana panah saka~$w$ menyang~$w'$. Kajaba iku,
kita menehi label simpul (donya) kanthi~$\mTrue{p}$ yen $w \in V(p)$,
lan yen ora kanthi~\mFalse{p}. \olref{fig:simple} nuduhake model
kanthi $W = \{w_1, w_2, w_3\}$, $R = \{\tuple{w_1, w_2},\tuple{w_1,
  w_3}\}$, $V(p) = \{w_1, w_2\}$, lan $V(q) = \{w_2\}$.

\begin{figure}
  \begin{center}
    \begin{tikzpicture}[modal]
      \node[world] (w1) [label={[align=right]right:\mTrue{p}\\ \mFalse{q}}]{$w_1$}; 
      \node[world] (w2) [label={[align=right]right:\mTrue{p}\\ \mTrue{q}},
        above right=of w1]{$w_2$}; 
      \node[world] (w3) [label={[align=right]right:\mFalse{p}\\ \mFalse{q}},
        below right=of w1] {$w_3$};
      \draw[->] (w1) to (w2);
      \draw[->] (w1) to (w3);
    \end{tikzpicture}
  \end{center}
  \caption{Model prasaja.}
  \ollabel{fig:simple}
\end{figure}

\end{document}
